\section{Use Cases per CTF Manager}\label{use-cases-per-ctf-manager}

Per rendere completa l\textquotesingle esposizione, sono elencati di
seguito gli Use Case per il CTF Manager, arricchito con alcuni dettagli
implementativi per comprendere al meglio l\textquotesingle architettura
del software.

Di seguito e\textquotesingle{} riportato un breve glossario:

\begin{quote}
\passthrough{\lstinline!br-team-k!}: e\textquotesingle{} il bridge di
rete che fornisce una rete locale per i container del team \(k\). Ogni
\passthrough{\lstinline!br-team-k!} e\textquotesingle{} logicamente
separato dagli altri \passthrough{\lstinline!br-team-k!}.
\end{quote}

\begin{quote}
\passthrough{\lstinline!wg-k!}: e\textquotesingle{}
l\textquotesingle interfaccia WireGuard che offre la connettivita VPN
per il team \(k\) verso \passthrough{\lstinline!br-team-k!}.
\end{quote}

\begin{quote}
\passthrough{\lstinline!eth0!}: in questo documento, spesso identifica
il nome dell\textquotesingle interfaccia che fornisce la connettivita
alla LAN per un container.
\end{quote}

\begin{quote}
\passthrough{\lstinline!verify-k!}: e\textquotesingle{} il bridge di
rete (gestito tramite SDN) utilizzato per implementare la rete di
verifica per il team \passthrough{\lstinline!k!}. Questa rete
e\textquotesingle{} completamente scollegata da
\passthrough{\lstinline!br-team-k!} e da \passthrough{\lstinline!wg-k!}.
\end{quote}

\begin{quote}
\passthrough{\lstinline!database!}: e\textquotesingle{} il database SQL,
usato da BlueAgent per tenere traccia delle risorse allocate.
\end{quote}

\subsection{\texorpdfstring{\texttt{UC-01}: Deploy
dell\textquotesingle infrastruttura di rete per il team
\(k\)}{UC-01: Deploy dell\textquotesingle infrastruttura di rete per il team k}}\label{uc-01-deploy-dellinfrastruttura-di-rete-per-il-team-k}

\begin{itemize}
\tightlist
\item
  \textbf{Attore}: CTF Manager
\item
  \textbf{Obiettivo}: Effettuare il deploy delle interfacce di rete
  necessarie al team \(k\) per eseguire i container e consentire
  l\textquotesingle accesso VPN ai partecipanti.
\item
  \textbf{Precondizione}: BlueAgent e\textquotesingle{} attivo ed
  e\textquotesingle{} raggiungibile via API REST.
\end{itemize}

Flusso delle azioni:

\begin{enumerate}
\def\labelenumi{\arabic{enumi}.}
\tightlist
\item
  Il CTF Manager invia la richiesta
  \passthrough{\lstinline!POST /teams!}, specificando
  l\textquotesingle identificativo \(k\) del team ed il relativo nome.
\item
  BlueAgent verifica se l\textquotesingle infrastruttura per il team
  \(k\) esiste gia:

  \begin{enumerate}
  \def\labelenumii{\arabic{enumii}.}
  \tightlist
  \item
    Se il team e\textquotesingle{} gia registrato, risponde
    \passthrough{\lstinline!HTTP 200 OK!}.
  \item
    Altrimenti, l\textquotesingle infrastruttura di rete viene creata da
    zero:

    \begin{enumerate}
    \def\labelenumiii{\arabic{enumiii}.}
    \tightlist
    \item
      Viene creato il bridge \passthrough{\lstinline!br-team-k!} per il
      team \(k\) con IP per il gateway
      \passthrough{\lstinline!10.k.1.1/24!}.
    \item
      Viene creata l\textquotesingle interfaccia WireGuard
      \passthrough{\lstinline!wg-k!} per il team \(k\). La porta UDP
      e\textquotesingle{} \passthrough{\lstinline!51820+k!},
      l\textquotesingle IP per il gateway e\textquotesingle{}
      \passthrough{\lstinline!10.k.0.1/24!}.
    \item
      Viene registrato il team all\textquotesingle interno del database.
    \item
      BlueAgent risponde con stato
      \passthrough{\lstinline!HTTP 201 Created!} confermando
      l\textquotesingle avvenuta creazione.
    \end{enumerate}
  \end{enumerate}
\end{enumerate}

Se avviene un errore irrimediabile durante la creazione, il bridge e
l\textquotesingle interfaccia WireGuard create vengono rilasciate.

\subsection{\texorpdfstring{\texttt{UC-02}: Deploy della challenge per
il team
\(k\)}{UC-02: Deploy della challenge per il team k}}\label{uc-02-deploy-della-challenge-per-il-team-k}

\begin{itemize}
\tightlist
\item
  \textbf{Attore}: CTF Manager
\item
  \textbf{Obiettivo}: Effettuare il deploy dei container necessari alla
  challenge da allestire per il team \(k\) richiedente.
\item
  \textbf{Precondizione}: \passthrough{\lstinline!UC-01!},
  \passthrough{\lstinline!UC-08!}. E\textquotesingle{} stato eseguito
  con successo il deploy dell\textquotesingle infrastruttura di rete per
  il team \(k\) e le immagini dei container richieste sono presenti nel
  sistema.
\end{itemize}

Flusso delle azioni:

\begin{enumerate}
\def\labelenumi{\arabic{enumi}.}
\tightlist
\item
  Il CTF Manager invia la richiesta
  \passthrough{\lstinline!POST /teams/\{team\_k\}/machines!}
  specificando \passthrough{\lstinline!team\_k!}.
\item
  BlueAgent valida le precondizioni.
\item
  Per ogni container specificato nella configurazione:

  \begin{enumerate}
  \def\labelenumii{\arabic{enumii}.}
  \tightlist
  \item
    Crea l\textquotesingle istanza del container.
  \item
    Collega l\textquotesingle interfaccia virtuale
    \passthrough{\lstinline!eth0!} dell\textquotesingle istanza al
    bridge del team \passthrough{\lstinline!br-team-k!}.
  \item
    Inietta nel filesystem dell\textquotesingle istanza la
    configurazione di rete statica.
  \item
    Avvia l\textquotesingle istanza del container ed applica
    l\textquotesingle indirizzo IP statico.
  \item
    Registra il container nel database.
  \end{enumerate}
\item
  BlueAgent restituisce la risposta con stato
  \passthrough{\lstinline!HTTP 201 Created!}.
\end{enumerate}

Se si verifica un errore durante l\textquotesingle avvio, BlueAgent
arresta ed elimina tutte le istanze e le risorse create fino a quel
momento per la challenge e risponde con un errore.

\subsection{\texorpdfstring{\texttt{UC-03}: Verifica della challenge per
il team
\(k\)}{UC-03: Verifica della challenge per il team k}}\label{uc-03-verifica-della-challenge-per-il-team-k}

\begin{itemize}
\tightlist
\item
  \textbf{Attore}: CTF Manager
\item
  \textbf{Obiettivo}: Verificare l\textquotesingle integrita dei servizi
  e l\textquotesingle efficacia delle patch applicate dal team mediante
  l\textquotesingle esecuzione dei comandi specificati nella
  configurazione della challenge.
\item
  \textbf{Precondizioni}: \passthrough{\lstinline!UC-02!}. I container
  relativi alla challenge specificata per il team \(k\) sono attivi.
\end{itemize}

Flusso delle azioni:

\begin{enumerate}
\def\labelenumi{\arabic{enumi}.}
\tightlist
\item
  Il CTF Manager invia la richiesta
  \passthrough{\lstinline!POST /teams/\{team\_k\}/verificator!}
  trasmettendo la configurazione della challenge da verificare.
\item
  Se una verifica per lo stesso team e\textquotesingle{} gia in corso,
  BlueAgent risponde con \passthrough{\lstinline!HTTP 409 Conflict!}.
\item
  Altrimenti, BlueAgent predispone l\textquotesingle ambiente di
  verifica:

  \begin{enumerate}
  \def\labelenumii{\arabic{enumii}.}
  \tightlist
  \item
    Crea uno switch virtuale denominato
    \passthrough{\lstinline!verify-k!}.
  \item
    Esegue la clonazione di tutti i container della challenge del team
    \(k\).
  \item
    Riconfigura l\textquotesingle interfaccia
    \passthrough{\lstinline!eth0!} dei cloni agganciandola allo switch
    di verifica.
  \item
    Avvia i cloni sandbox e configura al loro interno lo stesso
    indirizzo IP statico di produzione
    \passthrough{\lstinline!10.k.1.\{ip\_suffix\}/24!}.
  \item
    Se specificato nella configurazione, avvia un container di verifica.
  \end{enumerate}
\item
  Esecuzione dei controlli di verifica: viene eseguito un ciclo di
  tentativi finche\textquotesingle{} il comando non ha successo, oppure
  fino allo scadere del timeout.
\item
  Esegue i comandi di verifica all\textquotesingle interno del container
  specificato.
\item
  Dealloca le risorse sandbox create per il controllo.
\item
  BlueAgent risponde con stato \passthrough{\lstinline!HTTP 200 OK!}
  comunicando l\textquotesingle esito della verifica.
\end{enumerate}

Se si verifica un errore le risorse utilizzate vengono rilasciate e
viene comunicato l\textquotesingle errore.

\subsection{\texorpdfstring{\texttt{UC-06}: Dismissione e rilascio di
una challenge per il team
\(k\)}{UC-06: Dismissione e rilascio di una challenge per il team k}}\label{uc-06-dismissione-e-rilascio-di-una-challenge-per-il-team-k}

\begin{itemize}
\tightlist
\item
  \textbf{Attore}: CTF Manager
\item
  \textbf{Obiettivo}: Rimuovere tutti i container associati a una
  determinata challenge per il team \(k\).
\item
  \textbf{Precondizioni}: I container associati alla challenge
  specificata sono attivi per il team \(k\).
\end{itemize}

Flusso delle azioni:

\begin{enumerate}
\def\labelenumi{\arabic{enumi}.}
\tightlist
\item
  Il CTF Manager invia la richiesta
  \passthrough{\lstinline!DELETE /teams/\{team\_k\}/machines!}
  specificando nel payload la configurazione della challenge.
\item
  Per ciascun container specificato nella configurazione della
  challenge:

  \begin{enumerate}
  \def\labelenumii{\arabic{enumii}.}
  \tightlist
  \item
    Arresta l\textquotesingle istanza.
  \item
    Elimina l\textquotesingle istanza.
  \item
    Cancella il record del container dal database.
  \end{enumerate}
\item
  BlueAgent risponde con stato
  \passthrough{\lstinline!HTTP 204 No Content!}.
\end{enumerate}

In caso di problemi durante il teardown di uno o piu container,
BlueAgent restituisce un errore HTTP con i dettagli
dell\textquotesingle anomalia.

\subsection{\texorpdfstring{\texttt{UC-08}: Compilazione e Import di
un\textquotesingle immagine per
container}{UC-08: Compilazione e Import di un\textquotesingle immagine per container}}\label{uc-08-compilazione-e-import-di-unimmagine-per-container}

\begin{itemize}
\tightlist
\item
  \textbf{Attore}: CTF Manager
\item
  \textbf{Obiettivo}: Inviare la specifica di
  un\textquotesingle immagine in formato YAML ed avviare la compilazione
  e l\textquotesingle importazione nel catalogo delle immagini.
\item
  \textbf{Precondizione}: BlueAgent e\textquotesingle{} attivo.
\end{itemize}

Flusso delle azioni:

\begin{enumerate}
\def\labelenumi{\arabic{enumi}.}
\tightlist
\item
  Il CTF Manager invia una richiesta
  \passthrough{\lstinline!POST /images!} inviando una form multipart con
  il file di definizione YAML.
\item
  BlueAgent valida la specifica YAML e ne calcola l\textquotesingle hash
  crittografico.
\item
  BlueAgent verifica lo stato dell\textquotesingle immagine nel
  database:

  \begin{enumerate}
  \def\labelenumii{\arabic{enumii}.}
  \tightlist
  \item
    Se l\textquotesingle immagine con lo stesso alias o hash esiste gia,
    risponde immediatamente con \passthrough{\lstinline!HTTP 200 OK!}.
  \item
    Altrimenti, registra un nuovo record nel database.
  \end{enumerate}
\item
  BlueAgent avvia la compilazione tramite un worker asincrono:

  \begin{enumerate}
  \def\labelenumii{\arabic{enumii}.}
  \tightlist
  \item
    Importa l\textquotesingle immagine generata nel catalogo dei
    container.
  \item
    Assegna l\textquotesingle alias all\textquotesingle immagine
    importata e aggiorna il record nel database.
  \end{enumerate}
\item
  BlueAgent restituisce immediatamente al CTF Manager lo stato
  \passthrough{\lstinline!HTTP 202 Accepted!}.
\end{enumerate}

\subsection{\texorpdfstring{\texttt{UC-12}: Registrazione e
configurazione WireGuard per un membro del team
\(k\)}{UC-12: Registrazione e configurazione WireGuard per un membro del team k}}\label{uc-12-registrazione-e-configurazione-wireguard-per-un-membro-del-team-k}

\begin{itemize}
\tightlist
\item
  \textbf{Attore}: CTF Manager
\item
  \textbf{Obiettivo}: Registrare un giocatore come membro del team
  \(k\), allocare un indirizzo IP nella subnet VPN, generare le chiavi
  crittografiche WireGuard e restituire il file di configurazione
  \passthrough{\lstinline!.conf!}.
\item
  \textbf{Precondizioni}: BlueAgent e\textquotesingle{} attivo. Se
  l\textquotesingle infrastruttura di rete del team \(k\) non
  e\textquotesingle{} ancora attiva, viene allestita contestualmente.
\end{itemize}

Flusso delle azioni:

\begin{enumerate}
\def\labelenumi{\arabic{enumi}.}
\tightlist
\item
  Il CTF Manager invia la richiesta
  \passthrough{\lstinline!POST /teams/\{team\_k\}/members!} fornendo
  nome ed id dell\textquotesingle utente.
\item
  BlueAgent verifica se l\textquotesingle infrastruttura del team \(k\)
  e\textquotesingle{} attiva; se assente, esegue
  \passthrough{\lstinline!UC-01!}.
\item
  BlueAgent verifica se il membro e\textquotesingle{} gia registrato nel
  database per il team \(k\):

  \begin{enumerate}
  \def\labelenumii{\arabic{enumii}.}
  \tightlist
  \item
    Se il membro esiste gia, BlueAgent fornisce la configurazione
    rispondendo con \passthrough{\lstinline!HTTP 200 OK!}.
  \item
    Se il membro e\textquotesingle{} nuovo:

    \begin{enumerate}
    \def\labelenumiii{\arabic{enumiii}.}
    \tightlist
    \item
      Alloca un indirizzo IP VPN libero per il team. Se lo spazio degli
      indirizzi e\textquotesingle{} esaurito, risponde segnalando
      l\textquotesingle esaurimento degli indirizzi.
    \item
      Se non fornita, genera una nuova coppia di chiavi asimmetriche.
    \item
      Registra il nuovo peer sull\textquotesingle interfaccia
      \passthrough{\lstinline!wg-k!}.
    \item
      Salva il record del membro nel database.
    \end{enumerate}
  \end{enumerate}
\item
  BlueAgent risponde con stato
  \passthrough{\lstinline!HTTP 201 Created!} restituendo la
  configurazione WireGuard generata.
\end{enumerate}
